% Einheit 2 von 5 – Fläche zwischen zwei Graphen (bank/flaecheninhalt-durch-integration/e2.jsonl, zusammenbau v0.1)
%% TODO Zeitmarke Einheit 2: Marken-Zeile nicht lesbar
%% TODO Zweigzeile Teil 1: was hier gelernt wird (Satz in Schülersprache aus dem Einheitstitel) – Einheit 2
\einheitenkopf[e2]{Einheit 2 von 5 · Fläche zwischen zwei Graphen}
\zweigzeile{Abitur GK}
\setcounter{aufgabe}{10}

%% TODO Titel: Ich-kann-Satz fehlt in der Bank (Platzhalter: Kettenname) – Nr. 11
%% TODO Anweisung: ein Satz über der ersten Teilaufgabe fehlt in der Bank – Nr. 11
\begin{aufgabe}{Fläche zwischen zwei Graphen}
%% TODO \rechenplatz{n}: Schrittzahl des Grundfalls fehlt in der Bank (nur bei fester Schreibform, 2.3 b)
\begin{teile}
\teil Graphen von $f$ und $g$ im Bild, Intervall von $x = -2$ bis $x = 2$ – wer liegt oben? \kreuz{$f$ liegt oben}\\ \kreuz{$g$ liegt oben}\\ \kreuz{die obere Funktion wechselt}

\begin{ksys}[xmin=-3,xmax=4,ymin=-5,ymax=5,ablesen]\funktion{4-\x^2}{f} \funktion{\x^2-4}{g}\end{ksys}
\teil $f(x) = x^2 + 2$, $g(x) = x$ – Fläche zwischen den Graphen von $x = 0$ bis $x = 3$?
\teil $f(x) = x^2 - 2x$, $g(x) = x$ – Fläche, die die Graphen einschließen?
\teil $f(x) = x^2 - 4x + 5$ und die Gerade $y = 5$ – eingeschlossene Fläche?
\teil $f(x) = (x + 1) \cdot e^{x}$, $g(x) = e^{x}$; eine Stammfunktion von $x \cdot e^{x}$ ist $(x - 1) \cdot e^{x}$ – Fläche zwischen den Graphen von $x = 0$ bis $x = 1$?
\teil Eine Fläche liegt zwischen dem Graphen von $f(x) = 4 - 0{,}25x^2$ (von $-4$ bis $4$) und der x-Achse; ausgespart ist das Stück unter dem Graphen von $g(x) = 1 - x^2$ (von $-1$ bis $1$) – Inhalt?
\teil $f(x) = x^2$ mit den Tangenten in $x = -1$ und $x = 1$ – Fläche zwischen dem Graphen und beiden Tangenten?
\teil Für $k > 0$ schließen die Graphen von $f(x) = x^3$ und $g_k(x) = k^2 x$ eine Fläche aus zwei Flächenstücken ein. Untersuche die Aussage: Für $k > 2$ ist der Inhalt größer als $8$. \hfill (Abitur 2025 LK)
\end{teile}
\end{aufgabe}

%% TODO Titel: Ich-kann-Satz fehlt in der Bank (Platzhalter: Kettenname) – Nr. 12
%% TODO Anweisung: ein Satz über der ersten Teilaufgabe fehlt in der Bank – Nr. 12
\begin{aufgabe}{Fläche zwischen zwei Graphen – Fehler finden, Begründen, Anwendung}
\begin{teile}
\teil $f(x) = x^2 + 1$, $g(x) = x + 3$, Schnittstellen $-1$ und $2$ – wo ist der Fehler? \rechnung{A &= \text{Integral von } -1 \text{ bis } 2 \text{ über } (x^2 - x - 2) \\ &= -4{,}5}
\teil Beide Graphen liegen ganz unter der x-Achse. Warum liefert das Integral über obere minus untere Funktion trotzdem die Fläche zwischen ihnen?
\teil Die Ränder eines Teichs verlaufen entlang $f(x) = -0{,}2x^2 + 1{,}6x + 1$ und $g(x) = 1$ für $0 \le x \le 8$ (Angaben in m). Die Folie kostet $12$ € je Quadratmeter, für den Überlapp kommen $20\,\%$ Fläche dazu – Kosten?
\end{teile}
\end{aufgabe}
